\subsection{An explicit Jacobian class beyond the coordinate-unit symbols}

\begin{statement}{A.88}{Proved}
{The joint character has a nonzero Jacobian component}
Let $C,P_0,P_2,c,R$ be the smooth projective curve and points
of A.69, and choose the root coordinate of $P_2$ so that
$X_3(P_2)\in1+4\Z_2$.
Then $C$ has genus $33$.
The class
\[
 \mathcal A=\left(\frac{T-c}{R-c},2\right)\in\operatorname{Br}(C)
 \tag{A.88.1}
\]
has a nonzero image in $H^1(\Q,J)$, where $J=\operatorname{Jac}(C)$,
under the exact Hochschild--Serre comparison specified below.
An explicit divisor cocycle represents this image.

Moreover, there are no constant class $\gamma\in\operatorname{Br}(\Q)$
and finite-order characters
$\chi_i:G_\Q\to\Q/\Z$ for which
\[
             \mathcal A=\gamma+\sum_{i=1}^3(X_i,\chi_i)
                                   \quad\text{in }\operatorname{Br}\Q(C).
 \tag{A.88.2}
\]
Thus the completion argument using global roots of the old
coordinates cannot dispose of this class by expressing it as
a constant plus cyclic symbols in those coordinate roots.
This assertion specifies the coordinate-symbol subgroup;
it makes no assertion about symbols in other functions on $C$.
\end{statement}
\begin{proof}
We use the normalized coordinates of (A.69.1)--(A.69.3):
\[
 \begin{gathered}
 c_1=0,\quad c_2=1,\quad c_3=-9563/59024,\quad
 R=1/17,\quad c=-11/68,\quad a_i=R-c_i,\\
 T-c_i=a_iX_i^4,\qquad
 P_0=(R,1,1,1),\qquad
 P_2=(16/17,2,1/2,r),\quad r^4=4341/869.
 \end{gathered}
 \tag{A.88.3}
\]
The geometric cover has degree $64$ and inertia of order
four at each of $c_1,c_2,c_3,\infty$.
The tame Riemann--Hurwitz calculation gives
\[
                  2g-2=64\bigl(-2+4(1-1/4)\bigr)=64.
 \tag{A.88.4}
\]
Thus $g=33$ and the reduced infinity fibre $D_\infty$
has degree $16$.

\emph{The divisor and the exact cocycle.}
Put $L=\Q(\sqrt2)$ and let $\sigma$ be its nonidentity automorphism.
On the unramified fibre at $T=c$ put
$\eta=4X_1X_2X_3$.
The exact fibre relation (A.69.10) gives
$\eta^4=-1$ and hence
$(\eta+\eta^{-1})^2=2$.
Over $L$ define the effective divisors
\[
 Z_+:\ T=c,\quad\eta+\eta^{-1}=\sqrt2,\qquad
 Z_-:\ T=c,\quad\eta+\eta^{-1}=-\sqrt2.
 \tag{A.88.5}
\]
These are actual unions of geometric points in that etale
fibre, not divisor classes chosen by descent.
The $64$ root tuples give each of the four primitive
eighth roots $\eta$ with multiplicity $16$.
Each displayed sign selects two values of $\eta$,
so $\deg Z_+=\deg Z_-=32$.
The divisors are defined over $L$, and $\sigma Z_+=Z_-$.
For
$f=(T-c)/(R-c)$ one has
\[
 \operatorname{div}(f)=Z_++Z_- -4D_\infty,\qquad
 D=Z_+-2D_\infty,\qquad
                         D+\sigma D=\operatorname{div}(f).
 \tag{A.88.6}
\]
In particular $D$ has degree zero and
$\sigma[D]=-[D]$ in $J(L)$.
The map
\[
 z:G_\Q\longrightarrow J(\overline\Q),\qquad
 z(g)=
 \begin{cases}
 0,&g|_L=1,\\
 [D],&g|_L=\sigma
 \end{cases}
 \tag{A.88.7}
\]
is a continuous one-cocycle: its only nontrivial cocycle
identity is $[D]+\sigma[D]=0$.

Here is the comparison with the Brauer class, including
its sign convention.
The normalized cyclic two-cocycle representing the algebra
$(L(C)/\Q(C),f)$ is $f$ on $(\sigma,\sigma)$ and one
on the other pairs.
This algebra is the quaternion $(2,f)$, which is the same
quaternion class as $(f,2)$ by interchanging its generators.
For additive divisor cochains use
\[
              (\delta E)(g,h)=gE(h)-E(gh)+E(g).
 \tag{A.88.8}
\]
Take $E(g)=0$ or $D$ according as $g$ fixes $L$ or not.
Equation (A.88.6) says literally
$\delta E=\operatorname{div}(\text{the cyclic two-cocycle})$.
In the total-complex convention
$d_{\rm tot}=\delta+(-1)^p\operatorname{div}$
on group-cochain degree $p$, the pair consisting of
the cyclic two-cocycle and $-E$ is therefore a total
two-cocycle.  The edge map to
$H^1(\Q,\operatorname{Pic}\overline C)$ sends it to $-[z]$.
Since
\[
        \sigma[D]-[D]=-2[D],
 \tag{A.88.9}
\]
the cocycles $z$ and $-z$ differ by the coboundary of $[D]$.
Thus the edge image is $[z]$ as well; there is no
unresolved sign in its cohomology class.
This is the divisor-complex description of the
Hochschild--Serre map in
\cite[Proposition 4.3.4 and Remark 4.3.5]{AgentCTSkoro2019}.

For completeness the precise target here is the Jacobian.
Tsen's theorem gives
$\operatorname{Br}(\overline C)=0$, and the rational point
$P_0$ splits the structure morphism and the degree map
on the geometric Picard group.
The actual maps are
\[
 \begin{split}
 \ker\!\left(P_0^*:\operatorname{Br}(C)\to\operatorname{Br}(\Q)\right)
   &\xrightarrow{\ \epsilon\ }
           H^1(\Q,\operatorname{Pic}\overline C)
        \xrightarrow{\ \sim\ }H^1(\Q,J),\\
 \operatorname{Pic}\overline C&=
                   J(\overline\Q)\oplus\Z[P_0].
 \end{split}
 \tag{A.88.10}
\]
The first map is an isomorphism: the exact
Hochschild--Serre sequence has kernel the constant
Brauer group, and $P_0$ also makes its next
$H^3(\Q,\overline\Q^\times)\to H^3(C,\mathbb G_m)$ map
injective.  The second map is an isomorphism because
$H^1(\Q,\Z)=0$.
These exact sequences are
\cite[Proposition 4.3.2, Remark 4.3.3, equation (4.16),
Theorem 4.5.1(ii), and Remark 4.5.2]{AgentCTSkoro2019}.
The class $\mathcal A$ evaluates to zero at $P_0$, so its
image under (A.88.10) is exactly the cocycle (A.88.7).
Its evaluation at $P_2$ has invariant $1/2$ by A.69.
No constant class vanishing at $P_0$ can have this
evaluation, so $[z]\ne0$.

\emph{Residues restrict every putative coordinate character.}
Suppose (A.88.2) held.
At a boundary point over $T=c_i$, the function $X_i$
has order one and every $X_j$, $j\ne i$, is a unit.
Indeed $T-c_i=a_iX_i^4$ is the actual ramified local
parameter equation and the other punctures are distinct.
The residue of the right side is consequently
$\chi_i|_{G_{\kappa(P)}}$.
The left side is unramified by A.69.
Thus $\chi_i$ vanishes on the absolute Galois group of
every residue field at this boundary.
These residue fields are exactly the field factors of
\[
 \Q[Z_j^{\pm1}:j\ne i]/
                 (Z_j^4-(c_i-c_j)/a_j:j\ne i).
 \tag{A.88.11}
\]
The boundary constants are the following exact rational numbers:
\[
 \begin{array}{c|cc}
 i&\multicolumn{2}{c}{(c_i-c_j)/a_j,\quad j\ne i}\\ \hline
 1&17/16&9563/13035\\
 2&17&68587/13035\\
 3&-9563/3472&68587/55552 .
 \end{array}
 \tag{A.88.12}
\]
In the first two rows all constants are positive, so
A.80 makes the common abelian field exactly the
multiquadratic field generated by those two constants.
The last row also satisfies A.80's hypothesis.
Indeed
\[
 9563=73\cdot131,\quad 68587=107\cdot641,\quad
 3472=2^4\cdot7\cdot31,\quad
 55552=2^8\cdot7\cdot31 .
 \tag{A.88.13}
\]
The first squareclass has an odd valuation at $73$
and the second an odd valuation at $107$, with the
opposite class having zero valuation at that prime.
No nonempty product of them can be $-1$ modulo squares.
Therefore every $\chi_i$ in a putative (A.88.2) is
a sum of the two quadratic characters in its row.
This conclusion applies to arbitrary initial
finite-order characters; no order-two assumption
was imposed on them.

\emph{Every resulting coordinate symbol has equal evaluations.}
Choose $r\in1+4\Z_2$ as follows.
Since
\[
 \frac{4341}{869}=1+16\frac{217}{869},\qquad
 (1+4w)^4
      =1+16\bigl(w+6w^2+16w^3+16w^4\bigr),
 \tag{A.88.14}
\]
the polynomial inside parentheses has derivative one
modulo two and reduces to $w$.
Simple-root Hensel lifting gives a unique
$w\in\Z_2$ with its value $217/869$.
Then $r=1+4w$ is the required fourth root.
In fact $w$ is odd, so $r\equiv5\pmod8$.

Each of the four entries in the first two rows of
(A.88.12) is a square in $\Q_2$:
its valuation is even and its normalized unit is
one modulo eight.
Thus every symbol $(X_1(P_2),\chi_1)$ and
$(X_2(P_2),\chi_2)$ is zero.
The last row has even valuations $-4,-8$ and
normalized units five and three modulo eight.
Their products also have even valuation, so their
squareclasses have odd-unit representatives.

If $r\in1+4\Z_2$ and $b$ is any odd unit, then
$(r,b)=0$ in $\operatorname{Br}(\Q_2)$.
Here is a direct norm proof.
The squareclass of $r$ is one or five.
For the latter case let
$\omega=(1+\sqrt5)/2$.
Then
\[
                 N_{\Q_2(\sqrt5)/\Q_2}(a+\omega)
                                      =a^2+a-1.
 \tag{A.88.15}
\]
For every odd $b$ the equation
$a^2+a-1=b$ has a solution in $\Z_2$ by Hensel at
$a=0$: its value there is even and its derivative is
one modulo two.  Thus every odd unit is a norm,
which proves the asserted quaternion splitting.
For the squareclass one the assertion is immediate.
Applying this to the last row of (A.88.12) proves
$(X_3(P_2),\chi_3)=0$ as well.

All three coordinate symbols also vanish at $P_0$,
where every $X_i=1$.
The right side of (A.88.2) consequently has the same
$\Q_2$ evaluation at $P_0$ and at this chosen $P_2$,
whereas the left side has evaluations zero and
invariant $1/2$.
This contradiction proves (A.88.2) impossible.
Under the isomorphism (A.88.10), it says precisely
that the explicit Jacobian cocycle lies outside the
image of the unramified coordinate-symbol subgroup
modulo constants.
\end{proof}

This computation identifies the additional Picard obstruction
for one fixed old curve.  A.90 uses the individual coordinate
evaluations just proved to construct an ordinary completed old
tuple through the obstructed neighborhoods of A.85.
For that tuple the new coordinate is already obstructed modulo
two.  The exact failure of (A.88.2) is retained throughout
that construction.

The exact certificate \path{code/picard_obstruction_certificate.py}
checks the boundary constants and their squareclasses, the
finite divisor-cocycle identities and sign coboundary,
the fourth-root construction to high dyadic precision,
and every quadratic character combination forced by
(A.88.12).  Its output is retained beside the code.
